\section*{Capítulo \ref{ch:b2:asexual-reproduction} --- Reprodução assexuada e clonagem nas plantas}

\begin{solution}{exo:b2:asexual-reproduction:1}
\omterm{def:b2:asexual-reproduction:clone}{Clone}: os descendentes geneticamente idênticos de um único genitor,
obtidos por mitose. \omterm{def:b2:asexual-reproduction:clone}{Geneta}: o indivíduo genético, tudo o que cresceu a
partir de um zigoto. \omterm{def:b2:asexual-reproduction:clone}{Rameta}: uma unidade fisiologicamente independente
de uma \omterm{def:b2:asexual-reproduction:clone}{geneta}. Órgãos: estolões (morangueiro), rizomas
(grama-francesa), tubérculos (batata), bulbos (cebola), rebentos
(álamo), plântulas foliares (\emph{Kalanchoe}).
\end{solution}

\begin{solution}{exo:b2:asexual-reproduction:2}
Os câmbios vasculares do enxerto e do porta-enxerto precisam estar
alinhados e em contato; um calo fecha o corte e diferencia xilema e
floema através dele. O porta-enxerto fornece as raízes (água, minerais,
vigor, tolerância ao solo, controle do porte), o enxerto fornece o ramo
que frutifica (a variedade). Cada parte conserva seu próprio genoma e
as células nunca se fundem: dois \omterm{def:b2:meiosis-heredity:vocabulary}{genótipos} lado a lado, uma quimera;
um híbrido teria uma mistura dos dois genomas em cada célula.
\end{solution}

\begin{solution}{exo:b2:asexual-reproduction:3}
\omterm{def:b2:asexual-reproduction:clone}{Multiplicação vegetativa}: uma parte do corpo enraíza e forma uma nova
planta. \omterm{def:b2:asexual-reproduction:apomixis}{Apomixia}: uma \omterm{def:b2:angiosperm-reproduction:seed}{semente} se forma sem fecundação, e o embrião é
um \omterm{def:b2:asexual-reproduction:clone}{clone} da mãe. \omterm{def:b2:asexual-reproduction:apomixis}{Partenogênese}: um animal se desenvolve a partir de um
óvulo não fecundado. Só a \omterm{def:b2:asexual-reproduction:apomixis}{apomixia} conserva a \omterm{def:b2:angiosperm-reproduction:seed}{semente} --- com seu
envoltório, sua \omterm{def:b2:angiosperm-reproduction:seed}{dormência} e sua dispersão.
\end{solution}

\begin{solution}{exo:b2:asexual-reproduction:4}
Esterilize o explante; coloque-o em ágar com açúcar, minerais,
vitaminas e auxina e citocinina em equilíbrio (calo); aumente a razão
citocinina/auxina para induzir ramos; repique os ramos a cada poucas
semanas; transfira-os para um meio rico em auxina para que enraízem;
aclimate as plântulas sob umidade elevada e depois no solo.
\end{solution}

\begin{solution}{exo:b2:asexual-reproduction:5}
Cada planta é substituída por $1 + 24 = 25$: após 4 anos, $25^{4} =
\num{390625}$ plantas, que exigem \qty{39000}{m^2}, quase quatro
hectares. O espaço, a luz e os nutrientes, o curto alcance de um
estolão e a morte das plantas por adensamento encerram a fase
exponencial em um ou dois anos; o \omterm{def:b2:asexual-reproduction:clone}{clone} passa então a avançar como uma
frente.
\end{solution}

\begin{solution}{exo:b2:asexual-reproduction:6}
$\num{14000}/130 \approx 110$ gerações de troncos. Raio de um disco
de \qty{4.3e5}{m^2}: $\sqrt{4.3\times 10^{5}/\pi} = \qty{370}{m}$;
avanço de $370/\num{14000} = \qty{0.026}{m}$ por ano, menos de três
centímetros.
\end{solution}

\begin{solution}{exo:b2:asexual-reproduction:7}
$p_t = 2^{t}p_0/(1 + (2^{t} - 1)p_0)$: $t = 5$: $0.0032$; $t = 10$:
$0.1024/1.1023 = 0.093$; $t = 15$: $3.28/4.28 = 0.77$. Ele nunca se
espalha se sua aptidão for inferior à metade da aptidão da linhagem
sexuada: uma aptidão relativa $w$ dá $q' = 2wq$, que decai para $w <
\tfrac12$.
\end{solution}

\begin{solution}{exo:b2:asexual-reproduction:8}
$10 : 1$: raízes; $1 : 1$: calo indiferenciado; $1 : 10$: ramos. Uma
estaca já tem um ramo e precisa de raízes, que a auxina induz na base
cortada.
\end{solution}

\begin{solution}{exo:b2:asexual-reproduction:9}
1, 10, 100, \num{1000}, \num{10000} kg: após a terceira safra, apenas
\qty{1}{t}; após a quarta, \qty{10}{t}: quatro safras. A \omterm{def:b2:angiosperm-reproduction:seed}{semente}
verdadeira dá milhares de \omterm{def:b2:angiosperm-reproduction:seed}{sementes} por planta, mas a batata é um
tetraploide \omterm{def:b2:meiosis-heredity:vocabulary}{heterozigoto}, de modo que as plântulas segregam e nenhuma é
a variedade.
\end{solution}

\begin{solution}{exo:b2:asexual-reproduction:10}
$e^{-5} = 0.0067$: 670 indivíduos sem \omterm{def:b2:genome-diversification:mutation}{mutações} na população grande,
cerca de 7 na pequena. Sete indivíduos podem, todos, deixar de ter
descendentes por acaso em poucas gerações (a probabilidade de que
nenhum deles se reproduza é da ordem de $e^{-7}$ por geração), após o
que a melhor classe desaparece para sempre e a segunda melhor passa a
ser a melhor: um dente da catraca, repetido até que a carga se torne
insuportável. Com 670, a perda é praticamente impossível --- até que um
gargalo reduza $N$ ou a \omterm{thm:b2:genome-diversification:rate}{taxa de mutação} aumente; a catraca só gira num
sentido.
\end{solution}

\begin{solution}{exo:b2:asexual-reproduction:11}
No \omterm{def:b2:asexual-reproduction:clone}{clone}, todos os hospedeiros são suscetíveis: a linhagem se espalha
sem freio, à velocidade com que os esporos viajam. Na população
sexuada, o \omterm{def:b2:meiosis-heredity:vocabulary}{genótipo} que ela vence é uma de $2^{10} = 1024$
combinações, presente em cerca de um milésimo das plantas, espalhadas
entre vizinhas resistentes, de modo que a epidemia se esgota. A Lumper
era um único \omterm{def:b2:asexual-reproduction:clone}{clone} cultivado na maior parte de um país; a requeima
encontrou todas as plantas indefesas, e não havia variação para o
melhoramento até que outras variedades fossem importadas.
\end{solution}

\begin{solution}{exo:b2:asexual-reproduction:12}
A \omterm{prop:b2:asexual-reproduction:whysex}{catraca de Muller}, a combinação de \omterm{def:b2:genome-diversification:mutation}{mutações} favoráveis e a Rainha
Vermelha dão cada uma ao sexo um benefício que pode superar um fator
dois num mundo grande, parasitado e em mudança. Os rotíferos bdeloides
são o caso incômodo: quarenta milhões de anos sem machos e nenhum
colapso pela catraca. Eles parecem escapar por uma resistência extrema
à dessecação (que mata seus parasitas), pela captação de DNA estranho e
por um \omterm{prop:b2:genome-diversification:repair}{reparo do DNA} excepcionalmente eficiente --- exceções que
mostram o que a regra normalmente exige.
\end{solution}

\begin{solution}{pb:b2:asexual-reproduction:1}
\textbf{1.} $4\times 10^{4}$ plantas.
\textbf{2.} $1 + 6\times 2 = 13$.
\textbf{3.} \num{1300}; \num{16900}; \num{219700}.
\textbf{4.} $13^{t} = 400$: $t = \ln 400/\ln 13 = 5.99/2.56 = 2.3$
anos --- cheia durante o terceiro ano.
\textbf{5.} Raio de um disco de um hectare
$\sqrt{10^{4}/\pi} = \qty{56}{m}$; a \qty{0.5}{m} por ano, \num{113} anos.
\textbf{6.} Cem fundadoras espalhadas preenchem a lavoura
exponencialmente, a partir de cem frentes, em dois anos; uma única
fundadora preencheria sua vizinhança em dois anos e depois se
arrastaria por um século.
\textbf{7.} $5^{n} \ge 10^{6}$: $n \ge 8.6$, nove repicagens
($1.95\times 10^{6}$ ramos), 36~semanas.
\textbf{8.} $0.8\times 1.95\times 10^{6} = 1.6\times 10^{6}$ plantas.
\textbf{9.} $0.98$; $0.98^{10} = 0.82$.
\textbf{10.} $10\times 0.98 = 9.8$ linhagens limpas esperadas; a
partir de dez estacas, uma.
\textbf{11.} O vírus se desloca pelos plasmodesmos e pelo floema, e o
domo apical não tem conexão vascular e se divide mais depressa do que
o vírus se espalha, de modo que o décimo de milímetro mais externo
costuma estar livre de infecção.
\textbf{12.} Laboratório: plantas uniformes e livres de vírus de
imediato, mas caras, e qualquer \omterm{def:b2:genome-diversification:mutation}{mutação} ou contaminação na cultura é
multiplicada um milhão de vezes. Estolões: gratuitos e no próprio local,
mas com três safras perdidas e todos os vírus das plantas-mães levados
adiante.
\textbf{13.} Os apomíticos são $pN$ e dão $F$ \omterm{def:b2:angiosperm-reproduction:seed}{sementes} cada; os
sexuados, $(1 - p)N$, dão $F/2$ \omterm{def:b2:angiosperm-reproduction:seed}{sementes} cada (metade de seu esforço é
\omterm{def:b2:plant-life-cycles:heterospory}{pólen}, que não produz \omterm{def:b2:angiosperm-reproduction:seed}{sementes}); a parcela de \omterm{def:b2:angiosperm-reproduction:seed}{sementes} dos apomíticos é
$Fp/(Fp + F(1 - p)/2)
= 2p/(1 + p)$.
\textbf{14.} $p_5 = 0.032/1.031 = 0.031$; $p_{10} = 1.024/2.023 =
0.51$: ele passa de metade na geração 10.
\textbf{15.} $p' = 2p(1 - cp)/\bigl(2p(1 - cp) + (1 - p)\bigr)$.
\textbf{16.} $p' = p$ quando $2(1 - cp^*) = 1$: $p^* = 1/2c = 0.625$.
Abaixo de $p^*$, $2(1 - cp) > 1$ e $p$ aumenta; acima, diminui:
estável.
\textbf{17.} $p^* \ge 1$ quando $c \le \tfrac12$: o apomítico domina,
a menos que o parasita lhe custe mais da metade de sua produção quando
ele é comum.
\textbf{18.} $c = 0.6$: $p^* = 0.83$. Para manter o sexo na população,
o parasita precisa, com alta frequência de apomíticos, reduzir a
aptidão do \omterm{def:b2:asexual-reproduction:clone}{clone} em mais do que a vantagem dupla --- a Rainha Vermelha
precisa correr depressa.
\textbf{19.} Três conjuntos de cromossomos não conseguem parear na
\omterm{def:b2:meiosis-heredity:meiosis}{meiose}, de modo que seus gametas são desbalanceados; ele fica excluído
da recombinação, e a \omterm{prop:b2:asexual-reproduction:whysex}{catraca de Muller} e os parasitas agem sobre ele
sem oposição --- e é por isso que as linhagens apomíticas de
dente-de-leão são jovens e continuam sendo refeitas a partir de
ancestrais sexuados.
\textbf{20.} $10^{-9}\times 4.8\times 10^{8} = 0.48$ por divisão.
\textbf{21.} $2\times \num{14000}\times 20 = 5.6\times 10^{5}$
divisões; $0.48\times 5.6\times 10^{5} = 2.7\times 10^{5}$
\omterm{def:b2:genome-diversification:mutation}{mutações}.
\textbf{22.} $2.7\times 10^{5}/4.8\times 10^{8} = 5.6\times 10^{-4}$ do
genoma; cerca de \num{5400} diferenças codificantes.
\textbf{23.} A competição entre linhagens celulares no meristema (uma
linhagem mutante que se divide mais devagar é deslocada) e entre
\omterm{def:b2:asexual-reproduction:clone}{rametas} (um tronco com uma \omterm{def:b2:genome-diversification:mutation}{mutação} ruim produz menos rebentos). É mais
fraca porque a maioria das \omterm{def:b2:genome-diversification:mutation}{mutações} somáticas é heterozigota e
mascarada, e nada as expõe: nenhuma \omterm{def:b2:meiosis-heredity:meiosis}{meiose} as torna homozigotas e
nenhum estágio de zigoto \omterm{def:b2:unicellular-diversity:unicellular}{unicelular} testa o genoma inteiro numa só
célula.
\textbf{24.} Não é uniforme: seus troncos diferem em centenas de
milhares de sítios. Mas as diferenças são em sua maioria neutras e
mascaradas, de modo que o \omterm{def:b2:asexual-reproduction:clone}{clone} é, na prática, um único \omterm{def:b2:meiosis-heredity:vocabulary}{genótipo} --- um
mosaico de variantes somáticas de uma \omterm{def:b2:asexual-reproduction:clone}{geneta}.
\textbf{25.} Lavoura cheia após 2.3 anos; um meristema dá
$1.6\times 10^{6}$ plantas em nove meses; com $c = 0.8$, o apomítico
se estabiliza em $p^* = 0.625$, e os sexuados ficam com
\qty{37.5}{\%}.
\end{solution}
